% !TeX root = ../main.tex
% ---------------------------------------------------------------------------
% Excavation Design by Ground Type
% Scope extension added from the classical-textbook crosscheck.
% Maps to: Brady & Brown (2006) Ch. 7 "Excavation design in massive elastic
% rock", Ch. 8 "in stratified rock", Ch. 9 "in blocky rock";
% Hoek & Brown (1980) Ch. 7 "Underground excavation failure
% mechanisms".
% Rationale: the classical books organise excavation design by ground type,
% which is the structural idea the current book most clearly lacks. This
% chapter supplies the organising frame that Chapter~\ref{ch:empirical-support}
% then applies empirically.
% ---------------------------------------------------------------------------
\chapter{Excavation Design by Ground Type}
\label{ch:ground-type-design}


\section{Ground type as the organising principle}
\label{sec:ground-type-as-the-organising-principle}

\subsection{Ground type as the organising principle}
\label{sec:ground-type-principle}
\scopenote{Brady and Brown's central structural idea: massive, stratified and
blocky rock fail by different mechanisms and therefore require different
analyses. Classification systems describe quality; ground type describes
behaviour.}

\subsection{Behaviour classes and their diagnostic features}
\label{sec:behaviour-classes}
\scopenote{Massive, stratified, blocky, heavily fractured, weak and
high-stress ground, with the field evidence that distinguishes them and the
governing failure mode of each.}

\section{Massive rock: stress and failure zone}
\label{sec:massive-rock-stress-and-failure-zone}

\subsection{Massive elastic rock: boundary stresses}
\label{sec:massive-boundary-stress}
\scopenote{Stress concentration at the excavation boundary as a function of
shape and $k$-ratio; tensile zones in the crown and invert; the shape that
minimises the worst boundary stress for a given stress field.}

\subsection{Delineating the failure zone}
\label{sec:failure-zone}
\scopenote{Comparing boundary and near-boundary stress against a strength
criterion to map the extent of overstressed rock, and using that extent to
size reinforcement length.}

\subsection{The zone of influence of an excavation}
\label{sec:zone-of-influence}
\scopenote{The distance at which an opening ceases to disturb the stress field
appreciably, the conventional 5\,\% criterion, and its use in setting model
boundaries and excavation separations.}

\section{Excavation geometry and orientation}
\label{sec:excavation-geometry-and-orientation}

\subsection{Excavation shape and size effects}
\label{sec:shape-size-effects}
\scopenote{Aspect ratio, corner radius and sidewall height; why increasing
span increases both stress and the probability of intersecting an adverse
structure.}

\subsection{Excavation orientation}
\label{sec:excavation-orientation}
\scopenote{Orientation relative to the principal stress directions and to the
dominant discontinuity sets, and the resolution of conflict between the two.}

\subsection{Weakness planes in otherwise massive rock}
\label{sec:weakness-planes}
\scopenote{A single fault or persistent joint intersecting an otherwise
competent rock mass, and why it often governs a design that classification
would call good ground.}

\section{Stratified rock}
\label{sec:stratified-rock}

\subsection{Stratified rock: roof bed separation}
\label{sec:roof-separation}
\scopenote{Bedded and laminated roofs, separation between beds, the effective
thickness of the immediate roof, and the influence of bed thickness on
stability.}

\subsection{Voussoir beam analysis}
\label{sec:voussoir}
\scopenote{The linear arch model for a jointed or fractured roof beam:
compressive arch development, buckling and crushing limits, and the maximum
stable span. Related script: \code{arc-beam.py}.}

\subsection{Large-deflection roof behaviour}
\label{sec:large-deflection-roof}
\scopenote{Snap-through instability, progressive deflection under
time-dependent loading, and the difference between a roof that sags and a roof
that fails.}

\section{Support of laminated roofs}
\label{sec:laminated-roof-support}
\scopenote{Bolting to suspend, to build a composite beam, or to clamp beds;
selecting between these mechanisms from bed thickness and available anchorage.}

\section{Blocky rock}
\label{sec:blocky-rock}

\subsection{Blocky rock: from geometry to stability}
\label{sec:removability}
\scopenote{Block identification and removability are established in
Chapter~\ref{ch:structural-analysis}; this section takes the identified blocks
and asks whether they are stable. The distinction matters: geometry says a
block \emph{can} move, mechanics says whether it \emph{will}.}

\subsection{Key block theory in design}
\label{sec:key-blocks}
\scopenote{Using Goodman \& Shi's key block result as a design screen: if the
critical key block is stabilised, the blocks it keys are secured with it.
Applying this to set a systematic bolting pattern.}

\subsection{Wedge stability analysis}
\label{sec:wedges}
\scopenote{Limit equilibrium for the crown and sidewall wedges constructed in
Section~\ref{sec:tunnel-wedges}: driving weight against joint shear strength
and clamping stress, factor of safety conventions, and the sensitivity to
water pressure on the joint planes. Procedure in
Appendix~\ref{app:calculations}.}

\section{Sidewalls and weak rock}
\label{sec:sidewalls-and-weak-rock}

\subsection{Sidewall and high-wall design in blocky rock}
\label{sec:blocky-sidewalls}
\scopenote{Sliding and toppling of sidewall blocks, the reduced clamping
stress on a relaxed wall, and support selection for tall walls.}

\subsection{Heavily fractured and weak rock}
\label{sec:heavily-fractured}
\scopenote{Where block-based analysis loses meaning and equivalent continuum
treatment resumes; the transition criteria between the two frameworks.}

\section{Selecting the analysis route}
\label{sec:analysis-route}
\scopenote{A decision table mapping ground type to the appropriate empirical,
analytical and numerical methods, extending the matrix in
Chapter~\ref{ch:construction}.}
